% answers_short.tex —— 2004 年·天津卷·文科 · 简短答案速查
% 由 tools/make_short_answers.py 自 answers.tex 的 \gksolhead/\gkanswer 提取，
% 并与 2004/2004_数学真题_天津文_答案要点.md 逐题比对后生成，勿手改（改源 MD 后重跑）。

\begin{gknotice}
\textbf{使用说明}：本节只给最终答案便于快速核对，完整推导见解析版。答案为编写者独立推导，\textbf{非官方标准答案}。
\end{gknotice}

\gkpart{选择题}
\begin{center}\begin{tabular}{ccccccccccccc}
\toprule
题号 & 1 & 2 & 3 & 4 & 5 & 6 & 7 & 8 & 9 & 10 & 11 & 12 \\
\midrule
答案 & D & A & B & A & C & A & A & B & D & C & C & D \\
\bottomrule
\end{tabular}\end{center}

\gkpart{填空题}
\begin{itemize}[leftmargin=2.6em]
  \item[\textbf{13.}] 80
  \item[\textbf{14.}] -1
  \item[\textbf{15.}] $\left(-\infty, -\dfrac{13}{4}\right)$
  \item[\textbf{16.}] 36
\end{itemize}

\gkpart{解答题}
\begin{itemize}[leftmargin=2.6em]
  \item[\textbf{17.}] （1）$-\dfrac{1}{3}$；（2）$-\dfrac{5}{6}$
  \item[\textbf{18.}] （1）$\dfrac{1}{5}$；（2）$\dfrac{3}{5}$；（3）$\dfrac{4}{5}$
  \item[\textbf{19.}] （1）证明见解析；（2）$\dfrac{\sqrt{5}}{5}$
  \item[\textbf{20.}] （1）证明见解析；（2）$d = 2$，$a_n = 2n$
  \item[\textbf{21.}] （1）单调递增区间为 $(-\infty, -1)$ 和 $(1, +\infty)$，单调递减区间为 $(-1, 1)$；极大值为 $2$；（2）证明见解析
  \item[\textbf{22.}] （1）$\dfrac{x^2}{6} + \dfrac{y^2}{2} = 1$，$e = \dfrac{\sqrt{6}}{3}$；（2）$y = \pm \dfrac{\sqrt{5}}{5}(x - 3)$（或 $x \pm \sqrt{5}y - 3 = 0$）
\end{itemize}
